\section{Preliminaries} \label{s:prelim} 

\subsection{Notation} 
We use standard Vinogradov notation, thus $A\ll B$ means that $\lvert A\rvert \leq C \lvert B\rvert$, 
for an absolute constant $C$. By $A=B+O(\delta)$ we mean that $A-B \ll \delta$.  

For integers $N$, we set $[N] = \{1,2, \ldots, N\}$. 
For a finite set $F$ and $f \colon F\to \mathbb C$, we write 
\begin{equation}
    \Expectation_F f = \frac{1}{\lvert F\rvert} 
    \sum_{n\in F} f(n). 
\end{equation}
And we will frequently write $\Expectation_{[N]}=\Expectation_N$.  
Averages over an interval of indices are written $\Expectation_{M\leq n<N}$, 
and a limit $\lim_{N-M\to\infty}$ is a limit over all such intervals, as their length tends to infinity.  
Such averages, and the limits of them, are referred to as \emph{uniform}.  
We will use repeatedly the following elementary observation, in the form given in Corollary \ref{c:syndetic}: 
if a non-negative sequence $(a_n)$ satisfies $\liminf_{N-M\to\infty} \Expectation_{M\leq n<N} a_n>0$, 
then $\{n \colon a_n >0\}$ is syndetic.  

\subsection{Measure Preserving Systems} \label{sub:mps}

We recall standard notions around measure preserving systems.  
In particular, $(X,\mathcal A,m,T)$ 
is a \emph{measure preserving system} if $(X,\mathcal A,m)$ is a probability space, and $T \colon X\to X$ preserves $m$, in that $m(T^{-1} A)=m(A)$ for all measurable $A\in \mathcal A$.  
By $\Expectation_X f $ we mean $\int_X f \;dm$. 

\emph{Standing assumption.}  From here to the end of \S\ref{sec:compact}, unless stated otherwise, $(X,\mathcal A,m,T)$ is an invertible measure preserving system 
in which $X$ is a compact metric space and $\mathcal A$ is its Borel $\sigma$-algebra; in particular $m$ is a Radon measure.  
A system is \emph{ergodic} if its $T$-invariant sets have measure $0$ or $1$; this is assumed only where stated.  
Theorem \ref{t:CSZ} concerns arbitrary invertible systems, and Proposition \ref{p:reduction} in \S\ref{sub:reduction} shows that it suffices to treat 
ergodic systems satisfying the standing assumption.  





In different settings, the underlying space $(X,\mathcal A, m )$ remains the same, while we change the measure preserving operator from $T$, to some power $T^k$. 
In this case, we  just indicate the change in the operator.  
In other cases, we have  $T$ invariant sigma algebras, $\mathcal D \subset \mathcal E \subset \mathcal A$. 
In that case, we may refer to the system by its sigma algebra. 

The system is called \emph{weak mixing} if for all $A, B\in \mathcal A$, one has 
\begin{equation} \label{e:weaklyMixingDef}
    \lim_N \Expectation_N \lvert 
    m(A\cap T^{-n}B) - m(A)m(B) \rvert =0. 
\end{equation}
This condition for sets extends to functions. For $f\in L^2(X)$ with mean zero,  
\begin{equation} \label{wm-functions}
    \lim_N \Expectation_N  
    \lvert \langle f, T^n g \rangle \rvert =0, \qquad g \in L^2(X). 
\end{equation}
Equivalently, the spectral measure $\sigma _f$ of a mean zero $f \in L^2(X)$ has no atoms.  
Here, and below, for a unitary operator $U$ on a Hilbert space and a vector $v$, the \emph{spectral measure} $\sigma _v$ 
is the finite positive Borel measure on $\mathbb T = \mathbb R/\mathbb Z$ with 
\begin{equation} \label{e:spectral}
\langle U ^{n}v, v\rangle = \int_{\mathbb T} e(n \alpha ) \;d \sigma _v ( \alpha ), \qquad n \in \mathbb Z . 
\end{equation}
When no operator is named, $U$ is the Koopman operator $f \mapsto f\circ T$ of the system at hand.  
  






\bigskip 

A function $f\in L^2(X)$ is called \emph{compact} if the set $\overline{ \{T^n f \colon n\in \mathbb{Z}\} }$ is a compact subset of $L^2 (X)$.  The system is \emph{compact} if every $f\in L^2(X)$  is compact.  In this case $T$ has pure point spectrum: $L^2(X)$ is spanned by eigenfunctions of $T$.  
The eigenvalues need not be roots of unity; an irrational rotation of the circle is a compact system.  


\smallskip 
Let $\mathcal D \subset \mathcal A$ be a $T$ invariant 
sigma algebra. The system 
$(X,\mathcal D, m , T)$ is said to be a \emph{factor} of 
the original system $(X,\mathcal A,m ,T)$.  
(Here, there is a slight abuse of notation, in that we use the same notation for the set function $m$ restricted to the collection $\mathcal D$.) 
Factors are then naturally identified with their sigma algebras. 
And, given $\mathcal D\subset \mathcal E \subset \mathcal A$, 
we refer to $ \mathcal E$ as an extension of $\mathcal D$.  

With $\mathcal D\subset \mathcal E \subset \mathcal A$, we define a conditional inner product 
\begin{equation}
    \langle f,g \rangle _{\mathcal E\mid \mathcal D } 
    = \mathbb{E}_m  (f\overline g \mid \mathcal D) 
    \in L^2(\mathcal D) , \qquad 
    f,g\in L^2(\mathcal E). 
\end{equation}
This conditional inner product defines conditional  $L^2$ norms by 
$ \lVert f \rVert _{ \mathcal E \mid \mathcal  D }$. 
A conditional Cauchy--Schwarz and triangle inequality hold for this norm. 

\begin{definition}
    The Hilbert module $L^2(\mathcal E\mid \mathcal D)$ 
    are all those $f\in L^2(\mathcal E)$ for which the conditional norm  $\lVert f \rVert _{ \mathcal E \mid \mathcal  D } \in L^\infty (\mathcal D)$. 
\end{definition}


\begin{definition}
    A function $f\in L^2 (\mathcal E) $ is said to be \emph{conditionally weakly mixing} if  
    \begin{equation}
        \lim_N \Expectation_{N} 
        \lvert \langle T^nf,f\rangle _{\mathcal E\mid \mathcal D} 
        \rvert =0, \quad \textup{in $L^2(\mathcal D)$} . 
    \end{equation}
The extension $\mathcal E$ of $\mathcal D$ is said to be \emph{weakly mixing} if every $f\in L^2 (\mathcal E \mid \mathcal D) $ with $\mathbb E(f\mid \mathcal D)=0$  is conditionally weakly mixing. 
\end{definition}

We will use the definition of a weakly mixing extension only through the following consequence, 
which is phrased in terms of the relatively independent self-joining.  
As $X$ is a compact metric space, $m$ has a disintegration $m = \int m_x \;dm(x)$ over $\mathcal D$ \cite{MoreiraNotes}*{Theorem 7.2.1}, 
so that $\mathbb E(a\mid \mathcal D)(x) = \int a \;dm_x$ for $a \in L ^{\infty}(X)$ and almost every $x$.  
The \emph{relatively independent self-joining} of $m$ over $\mathcal D$ is the probability measure 
$\lambda \coloneqq \int m_x \otimes m_x \;dm(x)$ on $X \times X$; it appears, for instance, in the proof of \cite{MoreiraNotes}*{Proposition 8.6.1}.  
It has the two properties we need.  First, for $a, a' \in L ^{\infty}(X)$, 
\begin{equation} \label{e:joining}
\int_{X \times X} a \otimes a' \;d \lambda = \int_X \mathbb E(a\mid \mathcal D)\, \mathbb E(a'\mid \mathcal D)\;dm . 
\end{equation}
Second, $\lambda$ is invariant under $R \coloneqq T \times T$.  Indeed, as $\mathcal D$ is $T$-invariant, $\mathbb E(a\circ T\mid \mathcal D) = \mathbb E(a\mid \mathcal D)\circ T$ 
\cite{MoreiraNotes}*{Lemma 8.1.4}; so by \eqref{e:joining}, and as $T$ preserves $m$, 
$\lambda(R ^{-1}(B\times B')) = \int (\mathbb E(\mathbf 1_B\mid \mathcal D)\, \mathbb E(\mathbf 1_{B'}\mid \mathcal D))\circ T\;dm = \lambda(B\times B')$ 
for Borel sets $B, B'$.  Rectangles generate the Borel sets of $X \times X$ and are closed under intersection, so $\lambda \circ R ^{-1} = \lambda$.  

\begin{lemma} \label{l:atomless}
Let $f \in L ^{\infty}(\mathcal E)$ be conditionally weakly mixing, and set $G \coloneqq f \otimes \overline f$.  Then 
\begin{enumerate}
\item[(i)] the spectral measure $\sigma _f$ of $f$, for the Koopman operator of $T$ on $L^2(X)$, has no atoms; 
\item[(ii)] the spectral measure $\sigma _G$ of $G$, for the Koopman operator of $R$ on $L^2(\lambda)$, has no atoms, and 
\begin{equation} \label{e:psiIdentity}
\langle R ^{n}G, G\rangle _{L^2(\lambda)} = \bigl\lVert \langle T ^{n}f, f\rangle _{\mathcal E\mid \mathcal D}\bigr\rVert _{L^2(\mathcal D)} ^2 , 
\qquad n \in \mathbb Z . 
\end{equation}
\end{enumerate}
\end{lemma}

\begin{proof}
The identity \eqref{e:psiIdentity} is \eqref{e:joining} with $a = T ^{n}f\cdot \overline f$ and $a' = \overline{T ^{n}f}\cdot f = \overline a$, 
since $R ^{n}G\cdot \overline G = a \otimes a'$; the right side of \eqref{e:joining} is then $\int \lvert \mathbb E(T ^{n}f\cdot\overline f\mid \mathcal D)\rvert ^2 \;dm$.  

We use Wiener's lemma: for a finite positive measure $\sigma$ on $\mathbb T$, with $\widehat\sigma (n) = \int e(n \alpha )\;d \sigma (\alpha )$, 
\begin{equation}
\sum _{\theta} \sigma ( \{\theta\}) ^2 = \lim_N \Expectation_{1\leq n\leq N} \lvert \widehat\sigma (n)\rvert ^2 . 
\end{equation}

Both conclusions are unaffected by scaling $f$, so we may assume $\lVert f\rVert _{\infty} \leq 1$.  
Let $a_n \coloneqq \lVert \langle T ^{n}f,f\rangle _{\mathcal E\mid \mathcal D}\rVert _{L ^{1}(\mathcal D)}$.  
The functions $\lvert \langle T ^{n}f,f\rangle _{\mathcal E\mid \mathcal D}\rvert$ are non-negative and bounded by $1$, so 
\begin{equation}
\Expectation_{1\leq n\leq N} a_n = \Bigl\lVert \Expectation_{1\leq n\leq N} \bigl\lvert \langle T ^{n}f,f\rangle _{\mathcal E\mid \mathcal D}\bigr\rvert \Bigr\rVert _{L ^{1}(\mathcal D)} 
\leq \Bigl\lVert \Expectation_{1\leq n\leq N} \bigl\lvert \langle T ^{n}f,f\rangle _{\mathcal E\mid \mathcal D}\bigr\rvert \Bigr\rVert _{L ^{2}(\mathcal D)} \to 0 , 
\end{equation}
the limit being the definition of conditional weak mixing.  
Now $\widehat{\sigma _f}(n) = \langle T ^{n}f,f\rangle = \int \langle T ^{n}f,f\rangle _{\mathcal E\mid \mathcal D}\;dm$, so $\lvert \widehat{\sigma_f}(n)\rvert \leq a_n$; 
and $0 \leq \widehat{\sigma _G}(n) \leq a_n$ by \eqref{e:psiIdentity}, since the square of a function bounded by $1$ is at most its modulus.  
As $a_n \leq 1$, it follows that 
$\Expectation_{1\leq n\leq N} \lvert \widehat\sigma (n)\rvert ^2 \leq \Expectation_{1\leq n\leq N} a_n \to 0$ for $\sigma \in \{\sigma _f, \sigma _G\}$, 
and Wiener's lemma shows that neither measure has an atom.  
\end{proof}






\begin{definition}
    A function $f\in L^2 (\mathcal E \mid \mathcal D) $ is said to be \emph{conditionally compact} if for every $\epsilon>0$ there 
    are a finite number of functions $g_1 ,\ldots,g_r \in L^2(\mathcal E\mid \mathcal D)$ so that 
    \begin{equation}
        \sup_{n\in \mathbb Z} 
        \min_{1\leq q\leq r} 
        \bigl\lVert 
        \lVert T^n f - g_q \rVert_{\mathcal E\mid  \mathcal D} 
        \bigr\rVert_{L^\infty (\mathcal { D})} < \epsilon. 
    \end{equation}

    
    The extension $\mathcal E$ of $\mathcal D$ is said to be \emph{compact} if conditionally compact functions  are dense in  $L^2 (\mathcal E) $.  
\end{definition}








\subsection{Cantor Sets} 

We defined the Cantor sets of interest in \eqref{e:Cantor}. Here, we collect the properties of them needed to prove the main theorem.  
Throughout, $\mathcal C = \mathcal C_{b,D} = \{0=k_0<k_1 < \cdots\}$, with $0 \in D$, and we abbreviate 
\begin{equation} \label{e:delta}
\delta \coloneqq \lvert D\rvert \geq 2 . 
\end{equation}

The following identities describe the enumeration of $\mathcal C$; they are from \cite{firstpaper}*{Prop.~2.2}, 
and as everything below rests on them we include the short proof.  

\begin{lemma} \label{l:Diff}
    Let $\mathcal C=(k_n)_{n\geq 0}$ be an integer Cantor set in base $b$, with digit set $D=\{0=d_0<\cdots <d_{\delta -1}\}\subset \{0,\ldots ,b-1\}$. Then: 
    \begin{enumerate}
        \item[(i)] For any $n\in \mathbb{N}_0$, written in base $\delta $ as $n=\sum_{c=0}^tn_c\delta ^c$ with $n_c\in \{0,\ldots ,\delta -1\}$, we have $k_n=\sum_{c=0}^td_{n_c}b^c$.  
        \item[(ii)] For any $i\geq 0$, $n\geq 0$, and $0\leq j<\delta ^i$, we have $k_{\delta ^in+j}=b^ik_n+k_j$. 
    \end{enumerate}
\end{lemma}

\begin{proof}
For (i), define $\varphi \colon \mathbb N_0 \to \mathbb N_0$ by $\varphi(n) \coloneqq \sum_c d_{n_c} b ^{c}$, where the $n_c$ are the base $\delta$ digits of $n$.  
Since $d_0 = 0$, appending leading zero digits to $n$ does not change $\varphi(n)$, and the base $b$ digits of $\varphi (n)$ are exactly the $d_{n_c}$, all of which lie in $\{0 , \ldots, b-1\}$.  
So $\varphi(n) \in \mathcal C$; conversely, every element of $\mathcal C$ is $\varphi(n)$ for the $n$ whose base $\delta$ digits index its base $b$ digits in $D$, 
padding with the digit $0 = d_0$ where needed.  Thus $\varphi$ maps $\mathbb N_0$ onto $\mathcal C$.  It is strictly increasing, and so injective: if $n<n'$, let $c$ be the highest position at which their base $\delta$ digits differ, 
so that $n_c < n'_c$.  Then $\varphi(n)$ and $\varphi(n')$ have the same base $b$ digits above position $c$, and $d_{n_c} < d_{n'_c}$ at position $c$, so $\varphi(n)<\varphi(n')$.  
A strictly increasing map from $\mathbb N_0$ onto $\mathcal C$ is the increasing enumeration, so $k_n = \varphi(n)$.  

For (ii), write $j = \sum_{c<i} j_c \delta ^{c}$.  As $j< \delta ^{i}$, the base $\delta$ digits of $\delta ^{i}n + j$ are $j_0 , \ldots, j_{i-1}$ in the positions below $i$, 
followed by the digits of $n$, shifted up by $i$.  So by (i), $k_{\delta ^{i}n+j} = \sum_{c<i} d_{j_c}b ^{c} + \sum_c d_{n_c}b ^{c+i} = k_j + b ^{i}k_n$.  
\end{proof}


Note that part (ii) says that the blocks of indices 
$[\delta ^i n , \delta ^{i} (n+1))$, which partition $\mathbb N_0$ into consecutive intervals of length $\delta ^{i}$, 
correspond to the blocks $b ^{i} k_n + \bigl( \mathcal C \cap [0,b ^{i}) \bigr)$ of the Cantor set.  
This is the property that permits us to prove syndeticity, in place of mere positive density.  
It also controls the differences $k _{n+h}-k_n$, which are the subject of the van der Corput argument of \S\ref{sec:wm}.  

\begin{lemma} \label{l:differences}
Fix integers $h \geq 1$ and $\ell \geq 1$ with $\delta ^{\ell} > h$.  
For $0 \leq c < \delta ^{\ell}-h$ set $P_c \coloneqq \{ n \geq 0 \colon n \equiv c \mod \delta ^{\ell}\}$, 
and $\mathbb P \coloneqq \bigcup _{0\leq c < \delta ^{\ell}-h} P_c$. Then 
\begin{enumerate}
\item[(i)] For every $n \in P_c$, the difference $k _{n+h}-k_n$ is equal to the constant $\Delta _c \coloneqq k _{c+h}-k_c$.  
\item[(ii)] $\bigl\lvert \mathbb P \cap [M,N) \bigr\rvert \geq (1- h \delta ^{- \ell})(N-M) - \delta ^{\ell}$, for every interval $[M,N)$.  
\end{enumerate}
In particular, on a set of indices of uniform density at least $1-h \delta ^{-\ell}$, the difference $k_{n+h}-k_n$ takes at most $\delta ^{\ell}$ values, 
and each of the sets on which it is constant is a finite union of arithmetic progressions of step $\delta ^{\ell}$.  
\end{lemma}

\begin{proof}
For (i), write $n = \delta ^{\ell}a +c$ with $0 \leq c < \delta ^{\ell}-h$.  Then also $n+h = \delta ^{\ell}a + (c+h)$ with $0\leq c+h< \delta ^{\ell}$, 
so two applications of Lemma \ref{l:Diff}(ii) give $k_n = b ^{\ell}k_a + k_c$ and $k _{n+h} = b ^{\ell}k_a + k _{c+h}$.  Subtract.  
For (ii), $\mathbb P$ is the union of $\delta ^{\ell}-h$ residue classes modulo $\delta ^{\ell}$.  
\end{proof}


 

The distribution of Cantor sets in residue classes has been investigated in several papers, including \cites{ERDOS199899, SAAVEDRA-ARAYA_2026}.  
In \cite{firstpaper} it is shown that, for each $q \in \mathbb N$, the proportion of indices $n < N$ with $k_n \equiv a \mod q$ converges as $N \to \infty$, 
and that the limit is positive for $a = 0$.  The limiting distribution is in general not uniform on $\mathbb Z_q$, 
and the proportions do \emph{not} converge uniformly, in the sense of $\lim_{N-M\to\infty}$.  
Nevertheless, the set of indices $n$ with $q \mid k_n$ is not merely of positive density, but is syndetic.  
This is used only in Theorem \ref{t:compactSystem}, which is not needed for the main result; 
we give the short proof, as it is a good illustration of the block structure of Lemma \ref{l:Diff}(ii).  

\begin{lemma} \label{l:syndeticResidues}
Let $\mathcal C = \mathcal C_{b,D} = \{0=k_0<k_1 < \cdots\}$ be a Cantor set with $0 \in D$.  
For every integer $q \geq 1$, the set 
\begin{equation}
S_q \coloneqq \{ n \in \mathbb N_0 \colon k_n \equiv 0 \mod q \} 
\end{equation}
is syndetic.  
\end{lemma}

\begin{proof}
For $i \geq 1$ let $R_i \subset \mathbb Z_q$ be the set of residues $k_j \bmod q$ for $0 \leq j < \delta ^{i}$; 
by Lemma \ref{l:Diff}(i), 
\begin{equation}
R_i = \Bigl\{ \sum_{c=0}^{i-1} d_c b ^{c} \bmod q \colon d_0 , \ldots, d_{i-1} \in D \Bigr\} . 
\end{equation}
Since $0 \in D$, the sets $R_i$ increase with $i$, hence  they stabilize: there is an $i_0$ with $R_i = R \coloneqq \bigcup_{i} R_i$ for all $i \geq i_0$.  

Let $c_0 \geq 0$ and $\tau \geq 1$ be such that $b ^{c+\tau} \equiv b ^{c} \bmod q$ for all $c \geq c_0$, and set 
$Y= \{ d b ^{c} \bmod q \colon d \in D ,\ c \geq c_0\}$, a finite subset of $\mathbb Z_q$.  
Each element of $Y$ is of the form $db ^{c} \bmod q$ for infinitely many exponents $c$.  
Therefore any finite sum of elements of $Y$, repetitions allowed, can be realized as 
$\sum_{c} d_c b ^{c} \bmod q$ with the exponents $c$ \emph{distinct}, and so lies in $R$.  
The collection of all such sums is a non-empty finite sub-semigroup of the group $\mathbb Z_q$, and hence is the subgroup $H \subset \mathbb Z_q$ generated by $Y$.  
Thus $H \subset R$.  As $H$ is finite, there is an $i_1$ so that every element of $H$ is such a sum with all exponents less than $i_1$, whence $H \subset R_i$ for $i \geq i_1$.  

Now fix $i \geq \max (c_0 , i_0 , i_1)$ and let $n \in \mathbb N_0$ be arbitrary.  
Writing $k_n = \sum_{c} d_c b ^{c}$ with $d_c \in D$, we have $b ^{i} k_n = \sum_{c} d_c b ^{c+i}$, and all the exponents $c+i$ are at least $c_0$.  
Hence $b ^{i} k_n \bmod q \in H$, and therefore also $- b ^{i} k_n \bmod q \in H \subset R_i$.  
Select $0 \leq j < \delta ^{i}$ with $k_j \equiv - b ^{i} k_n \bmod q$.  By Lemma \ref{l:Diff}(ii), 
\begin{equation}
k _{\delta ^{i} n+j} = b ^{i} k_n + k_j \equiv 0 \mod q . 
\end{equation}
That is, $S_q$ meets the block of indices $[\delta ^{i} n , \delta ^{i}(n+1))$, for every $n \geq 0$.  
As these blocks partition $\mathbb N_0$ into consecutive intervals of length $\delta ^{i}$, the set $S_q$ is syndetic, with gaps at most $2 \delta ^{i}$.  
\end{proof}


The input from our first paper is the following uniform exponential sum estimate.  

\begin{lemma} \label{l:irrational}
Let $\mathcal C = \mathcal C_{b,D} = \{0=k_0<k_1 < \cdots\}$ be a Cantor set with $0\in D$, and let 
$E_0 \subset \mathbb T$ be the set of $\beta$ for which $\gcd(D)\, \beta$ is a $b$-adic rational modulo $1$; it is a countable set of rationals.  
Then for every $\beta \notin E_0$, and in particular for every irrational $\beta$, 
\begin{equation} \label{e:irrationalUniform}
\lim_{L\to\infty} \sup_{M \geq 0} \Bigl\lvert \Expectation_{M\leq n<M+L} e ( \beta k_n) \Bigr\rvert =0 . 
\end{equation}
\end{lemma}

This is the case of the uniform estimate in \cite{firstpaper}*{\S 4} in which the frequency attached there to the sum of digits is $0$; 
in the notation of that paper, $( \beta ,0)$ is exceptional exactly when $\gcd(D)\beta$ is a $b$-adic rational.  
The supremum over the starting index $M$ is the essential point, and the reason for the word `uniform' throughout this paper.  
As all of \S\ref{sec:wm} rests on it, we include the proof, which in this case is short.  

\begin{proof}
Write $s \coloneqq \gcd(D)$, and recall $\delta = \lvert D\rvert$.  

\emph{Reduction to $s=1$.}  The digit set $D' \coloneqq \{d/s \colon d \in D\}$ again satisfies $0 \in D' \subset \{0 , \ldots, b-1\}$ and $2 \leq \lvert D'\rvert < b$, and $\gcd(D')=1$.  
The elements of $\mathcal C$ are $k_n = s k'_n$, where $k'_0 < k'_1 < \cdots$ are the elements of $\mathcal C_{b,D'}$, so that $e(\beta k_n) = e(s\beta\, k'_n)$; 
and $s\beta$ is a $b$-adic rational modulo $1$ exactly when $\beta \in E_0$.  
So it suffices to treat the case $s=1$, with $s\beta$ in place of $\beta$.  We assume from now on that $s=1$, so that $E_0$ is the set of $b$-adic rationals modulo $1$.  

\emph{The digit average.}  Let $\rho (\theta) \coloneqq \Expectation_{d \in D} e(d\theta)$.  Then $\lvert \rho (\theta)\rvert \leq 1$, with equality only if $\theta \in \mathbb Z$.  
Indeed, equality forces the numbers $e(d\theta)$, $d \in D$, to coincide, hence to equal $1$, as $0 \in D$; 
so $d\theta \in \mathbb Z$ for every $d \in D$, and then $\theta \in \mathbb Z$, as $\gcd(D)=1$.  

\emph{Products over blocks.}  By Lemma \ref{l:Diff}(i), as $j$ runs over $0 \leq j < \delta ^{\ell}$, the base $b$ digits of $k_j$ in the positions $0 , \ldots, \ell-1$ 
run independently over $D$, and $k_j$ has no other non-zero digits.  Hence 
\begin{equation} \label{e:Pell}
P_\ell (\beta) \coloneqq \Expectation_{0\leq j<\delta ^{\ell}} e(\beta k_j) = \prod_{c=0}^{\ell-1} \rho (b ^{c}\beta) . 
\end{equation}

\emph{Window split.}  Fix $i \geq 1$ and an interval $[M,N)$, and let $M_i \coloneqq \lceil M/\delta ^{i}\rceil$ and $N_i \coloneqq \lfloor N/\delta ^{i}\rfloor$.  
At most $2\delta ^{i}$ indices of $[M,N)$ lie outside $[\delta ^{i}M_i, \delta ^{i}N_i)$; when $N_i \leq M_i$ this interval is empty, and $N-M < 2\delta ^{i}$.  
The remaining indices are partitioned into the blocks $\{\delta ^{i}m + j \colon 0 \leq j < \delta ^{i}\}$, $M_i \leq m < N_i$.  
By Lemma \ref{l:Diff}(ii), $e(\beta k_{\delta ^{i}m+j}) = e(\beta b ^{i}k_m)\, e(\beta k_j)$, so the sum over a block is $\delta ^{i}P_i(\beta)\, e(\beta b ^{i}k_m)$, of modulus at most $\delta ^{i}\lvert P_i(\beta)\rvert$.  
There are $N_i - M_i \leq (N-M)/\delta ^{i}$ blocks, and therefore 
\begin{equation} \label{e:windowSplitE}
\Bigl\lvert \Expectation_{M\leq n<N} e(\beta k_n) \Bigr\rvert \leq \lvert P_i(\beta)\rvert + \frac{2\delta ^{i}}{N-M} , \qquad M \geq 0 . 
\end{equation}
The bound does not depend on $M$; this is where the uniformity comes from.  

\emph{The limit of $P_\ell$.}  Write $\lVert \theta\rVert$ for the distance from $\theta$ to the nearest integer.  
Suppose first that $\lVert b ^{c}\beta\rVert \geq \eta$ for infinitely many $c$, for some $\eta>0$.  
As $\rho$ is continuous and $1$-periodic, and $\lvert \rho\rvert <1$ off the integers, there is a $\kappa>0$ with $\lvert \rho (\theta)\rvert \leq 1-\kappa$ whenever $\lVert \theta\rVert \geq \eta$.  
So infinitely many factors in \eqref{e:Pell} have modulus at most $1-\kappa$, all have modulus at most $1$, and $P_\ell(\beta) \to 0$.  

Otherwise $\lVert b ^{c}\beta\rVert \to 0$, and we show that $\beta \in E_0$.  Choose $c_0$ with $\lVert b ^{c}\beta\rVert < 1/(2b)$ for all $c \geq c_0$.  
If $0 < \lVert b ^{c}\beta\rVert < 1/(2b)$, write $b ^{c}\beta = z + \theta$ with $z \in \mathbb Z$ and $\lvert \theta\rvert = \lVert b ^{c}\beta\rVert$; 
then $b ^{c+1}\beta = bz + b\theta$ with $0<\lvert b\theta\rvert < 1/2$, so that $\lVert b ^{c+1}\beta\rVert = b \lVert b ^{c}\beta\rVert$.  
Hence if $\lVert b ^{c_0}\beta\rVert \neq 0$, the quantities $\lVert b ^{c_0+r}\beta\rVert = b ^{r}\lVert b ^{c_0}\beta\rVert$ would eventually exceed $1/(2b)$, a contradiction.  
So $b ^{c_0}\beta \in \mathbb Z$, that is, $\beta$ is a $b$-adic rational.  

\emph{Conclusion.}  Let $\beta \notin E_0$ and $\epsilon > 0$.  By the previous step, $P_i(\beta) \to 0$, so we may fix $i$ with $\lvert P_i(\beta)\rvert < \epsilon$.  
By \eqref{e:windowSplitE}, $\sup_{M\geq 0}\lvert \Expectation_{M\leq n<M+L} e(\beta k_n)\rvert \leq \epsilon + 2\delta ^{i}/L$, which is less than $2\epsilon$ once $L > 2\delta ^{i}/\epsilon$.  
As $\epsilon$ was arbitrary, this proves \eqref{e:irrationalUniform}.  
\end{proof}



The form in which Lemma \ref{l:irrational} enters the proof is the following.  For $\beta \in \mathbb R$ and $L \geq 1$ set 
\begin{equation} \label{e:Phi}
\Phi _L (\beta) \coloneqq \sup _{M\geq 0} \Bigl\lvert \Expectation_{M\leq n<M+L} e ( \beta k_n) \Bigr\rvert . 
\end{equation}

\begin{lemma}[Uniform spectral lemma] \label{l:uniformSpectral}
Let $\sigma$ be a finite positive Borel measure on $\mathbb T$ without atoms, and let $s \geq 1$ be an integer.  Then 
\begin{equation}
\lim_{L\to\infty} \int_{\mathbb T} \Phi _L (s \alpha ) \;d \sigma ( \alpha ) = 0 . 
\end{equation}
\end{lemma}

\begin{proof}
Each $\Phi_L$ is a supremum of countably many continuous functions, hence Borel, and $0 \leq \Phi _L \leq 1$.  
By Lemma \ref{l:irrational}, $\Phi_L(s \alpha ) \to 0$ unless $s \alpha \in E_0$ modulo $1$.  
Each point of $E_0$ has $s$ preimages under $\alpha \mapsto s\alpha$, so these $\alpha$ form a countable set, which is $\sigma$-null as $\sigma$ has no atoms.  
Dominated convergence completes the proof.  
\end{proof}

For instance, if $f \in L^2(X)$ has an atomless spectral measure, then, by \eqref{e:spectral}, 
\begin{equation} \label{e:Uniform}
\Bigl\lVert \Expectation _{M\leq n<N} T ^{s k_n}f \Bigr\rVert _2 ^2 
= \int_{\mathbb T} \Bigl\lvert \Expectation _{M\leq n<N} e(s\alpha k_n) \Bigr\rvert ^2 \;d \sigma _f (\alpha) 
\leq \int_{\mathbb T} \Phi _{N-M}(s \alpha )\;d \sigma _f ( \alpha ) , 
\end{equation}
which tends to zero as $N-M\to\infty$, uniformly in $M$.  This is the uniform ergodic theorem along $\mathcal C$, for such $f$.  

\subsection{Structure of the proof} 

The method of proof follows the lines of ergodic proofs of Szemer\'edi's Theorem. To summarize this approach, we say that a measure preserving system $(X,\mathcal A,m,T)$ is \emph{Cantor-Szemer\'edi},  or \emph{CSZ}, 
if this holds: 
For all measurable $A\subset X$, with $m(A) >0$, all integers $t \geq 1$, all integers $i \geq 0$, and all Cantor sets $\mathcal C = \mathcal C_{b,D}$, with base $b$, 
digit set $D$, and $0\in D$, there holds 
\begin{equation}  \label{e:CSZdef}
    \liminf_{N-M\to\infty} \Expectation_{M\leq n<N} 
    \int_X  \prod _{s=0} ^t T^{s b ^{i} k_n}\mathbf 1_A \;dm >0. 
\end{equation}
Here, $\mathcal C = \{0=k_0<k_1<\cdots\}$.  
The averages above are uniform in the index $n$; by Corollary \ref{c:syndetic}, 
this implies that the set of indices $n$ which contribute is syndetic, which is the point of insisting on the uniformity.  

It is equivalent to ask for \eqref{e:CSZdef} with $\mathbf 1_A$ replaced by an arbitrary non-negative $f \in L ^{\infty}(X)$ 
which is not identically zero.  One direction is trivial.  For the other, choose $c>0$ with $m(A)>0$ for $A \coloneqq \{f>c\}$; 
then $f \geq c \mathbf 1_A$, and all the factors are non-negative, so the average for $f$ is at least $c ^{t+1}$ times the average for $\mathbf 1_A$.  
In particular, the case $i=0$ of \eqref{e:CSZdef} is exactly the conclusion \eqref{e:CSZ} of Theorem \ref{t:CSZ}.  

Two comments on the dilation parameter $i$ are in order.  First, the case $i=0$ is the statement of Theorem \ref{t:CSZ}, 
and the other cases are statements of the same kind: since $b ^{i} k_n = k _{\delta ^{i} n}$ by Lemma \ref{l:Diff}(ii), 
the condition for a given $i$ is \eqref{e:CSZdef} with $i=0$, with the index restricted to the progression $\delta ^{i} \mathbb N_0$.  
Put another way, for a \emph{fixed} Cantor set $\mathcal C$ of base $b$, the condition for the dilation $i$ is the condition for $i=0$, 
for the system $(X,\mathcal A,m,T ^{b ^{i}})$ and the same $\mathcal C$.  
We stress that this is a statement about one Cantor set at a time.  It does not say that a CSZ system $(X,\mathcal A,m,T)$ gives rise to 
CSZ systems $(X,\mathcal A, m, T ^{b ^{i}})$: the CSZ property quantifies over all Cantor sets, and for a second Cantor set the available 
dilations are powers of \emph{its} base, not of $b$.  
Second, carrying the parameter $i$ is not a luxury: the proof for compact extensions in \S\ref{sec:compact} 
uses the hypothesis on the factor with a dilation determined by that proof, and returns the conclusion 
for the extension with the dilation it was given.  It is the self-similarity $b ^{i}\mathcal C \subset \mathcal C$ of the Cantor set which makes this  work.  

We prove 

\begin{theorem} \label{t:ECSZ}
    Every ergodic system satisfying the standing assumption of \S\ref{sub:mps} is CSZ.  
\end{theorem}

This proves Theorem \ref{t:CSZ}.  By Proposition \ref{p:reduction}, Theorem \ref{t:ECSZ} implies that every invertible measure preserving system is CSZ; 
the case $i=0$, with the passage from sets to non-negative functions recorded after \eqref{e:CSZdef}, is Theorem \ref{t:CSZ}.  
Theorem \ref{t:CSZ-Integers} then follows by the Furstenberg correspondence principle, 
as explained in the introduction.  

The proof of Theorem \ref{t:ECSZ} follows the maximal-factor strategy of Furstenberg, Katznelson and Ornstein \cite{MR670131}, 
in the form presented by Moreira \cite{MoreiraNotes}*{\S 8.1}.  
A \emph{factor} of $(X,\mathcal A,m,T)$ is a $T$-invariant sub-$\sigma$-algebra $\mathcal D \subset \mathcal A$; it is CSZ if $(X,\mathcal D,m,T)$ is; 
and factors which agree up to null sets are identified.  
The proof rests on three statements: the CSZ property passes to weakly mixing and to compact extensions (Lemma \ref{l:Extensions}); 
the CSZ factors are closed under suprema of chains (Lemma \ref{l:chain}); and every proper factor has a non-trivial extension of one of these two kinds 
(Lemma \ref{l:dichotomy}).  We first record what is used about powers of $T$.  

\begin{lemma} \label{l:powers}
Let $k \geq 1$ be an integer.  
\begin{enumerate}
\item If $\mathcal E$ is a compact extension of $\mathcal D$ for $T$, then it is one for $T ^{k}$. 
\item Let $f,g \in L^2(\mathcal E\mid \mathcal D)$ be such that 
\begin{equation} \label{e:wmHypothesis}
\lim_{N} \bigl\lVert \Expectation_N \lvert \langle T ^{n}f,g \rangle _{\mathcal E \mid \mathcal D} \rvert \bigr\rVert _{L^2(\mathcal D)} =0 . 
\end{equation}
Then the same holds along the progression $k \mathbb N$, that is 
\begin{equation} \label{e:wmProgression}
\lim_{N} \bigl\lVert \Expectation_N \lvert \langle T ^{kn}f,g \rangle _{\mathcal E \mid \mathcal D} \rvert \bigr\rVert _{L^2(\mathcal D)} =0 . 
\end{equation}
In particular, if every $f \in L ^{\infty}(\mathcal E)$ with $\mathbb E(f\mid \mathcal D)=0$ is conditionally weakly mixing for $T$, 
then the same holds for $T ^{k}$. 
\end{enumerate}
\end{lemma}

\begin{proof}
For (1), the orbit $\{T ^{kn}f \colon n \in \mathbb Z\}$ is a subset of $\{T^nf \colon n \in \mathbb Z\}$, so conditional compactness is inherited, 
and the conditionally compact functions therefore remain dense in $L^2(\mathcal E)$. 

For (2), set $\phi_n \coloneqq \lvert \langle T^nf,g\rangle _{\mathcal E\mid \mathcal D}\rvert$ and $a_n \coloneqq \lVert \phi _n\rVert _{L^2(\mathcal D)}$.  
By the conditional Cauchy--Schwarz inequality, 
$0 \leq \phi _n \leq \lVert T ^{n} f\rVert _{\mathcal E\mid \mathcal D} \lVert g\rVert _{\mathcal E\mid \mathcal D} \leq C$ for a constant $C$, 
since $f, g$ lie in the Hilbert module.  
As the $\phi_n$ are non-negative, \eqref{e:wmHypothesis} gives 
\begin{equation}
\Expectation_N \lVert \phi _n \rVert_{L^1(\mathcal D)} = \bigl\lVert \Expectation_N \phi _n \bigr\rVert _{L^1(\mathcal D)} 
\leq \bigl\lVert \Expectation_N \phi _n \bigr\rVert _{L^2(\mathcal D)} \to 0 , 
\end{equation}
and then $\Expectation_N a_n ^2 \leq C \,\Expectation_N \lVert \phi_n\rVert _{L^1(\mathcal D)} \to 0$, 
so that $\Expectation_N a_n \to 0$ by Cauchy--Schwarz.  
A bounded non-negative sequence has $\Expectation_N a_n \to 0$ if and only if $a_n \to 0$ along a set $S$ of density one.  
The complement of $S$ has density zero, so $\lvert S ^{c} \cap [1, kN]\rvert = o(N)$ for fixed $k$; 
that is, $a _{kn} \to 0$ along a set of $n$ of density one, so that $\Expectation_N a _{kn} \to 0$.  
Since $\lVert \Expectation_N \phi_{kn}\rVert _{L^2(\mathcal D)} \leq \Expectation_N a _{kn}$, this proves \eqref{e:wmProgression}.  
For the last statement, take $g=f$ with $\mathbb E (f\mid \mathcal D)=0$: then \eqref{e:wmHypothesis} is the definition of 
$f$ being conditionally weakly mixing, and \eqref{e:wmProgression} is that same statement for $T^k$. 
\end{proof}

\begin{lemma}[Extension principles] \label{l:Extensions}
Let $\mathcal D \subset \mathcal E$ be factors of $(X,\mathcal A,m,T)$, and suppose that $\mathcal D$ is CSZ.  
\begin{enumerate}
\item[(W)] If $\mathcal E$ is a weakly mixing extension of $\mathcal D$, then $\mathcal E$ is CSZ. 
\item[(K)] If $\mathcal E$ is a compact extension of $\mathcal D$, then $\mathcal E$ is CSZ. 
\end{enumerate}
\end{lemma}

Part (W) is proved in \S\ref{sec:wm}, as a consequence of Lemma \ref{l:wmExtension}; part (K) is Theorem \ref{t:compactExtension}.  
The trivial factor, consisting of the null and the conull sets, is CSZ: if $A$ is in it and $m(A)>0$, then $A$ is conull, 
and every integral in \eqref{e:CSZdef} equals $1$.  
A weakly mixing system is a weakly mixing extension of the trivial factor, and a compact system is a compact extension of it, 
so Lemma \ref{l:Extensions} shows in particular that weakly mixing systems and compact systems are CSZ; see Corollary \ref{t:wm}.  
Theorem \ref{t:compactSystem}, which treats a restricted class of compact systems by a direct argument, 
is not used in the proof of the main theorem.  

\begin{lemma}[Chains] \label{l:chain}
Let $\{\mathcal D_\alpha\}_{\alpha \in I}$ be a non-empty family of CSZ factors of $(X,\mathcal A,m,T)$ which is totally ordered by inclusion, 
and let $\mathcal D$ be the $\sigma$-algebra generated by their union.  Then $\mathcal D$ is a CSZ factor.  
\end{lemma}

\begin{proof}
The $\sigma$-algebra $\mathcal D$ is $T$-invariant, as each $\mathcal D_\alpha$ is.  

\emph{Approximation.}  For every $E \in \mathcal D$ and $\eta>0$ there is an $\alpha$ with 
$\lVert \mathbf 1_E - \mathbb E(\mathbf 1_E \mid \mathcal D_\alpha)\rVert_2 < \eta$; see \cite{MoreiraNotes}*{Lemma 8.1.5}.  
For completeness: the union $\bigcup_\alpha \mathcal D_\alpha$ is an algebra of sets, since any two of its members lie in a common $\mathcal D_\alpha$; 
and $\bigcup _\alpha L^2(\mathcal D_\alpha)$ is a linear subspace of $L^2(X)$, with closure $\mathcal H$, say.  
The sets $B \in \mathcal D$ with $\mathbf 1_B \in \mathcal H$ contain that algebra, and they form a monotone class, since monotone limits of indicators converge in $L^2$.  
By the monotone class theorem, they are all of $\mathcal D$.  So for some $\alpha$ there is a $g' \in L^2(\mathcal D_\alpha)$ with $\lVert \mathbf 1_E - g'\rVert_2 < \eta$, 
and $\mathbb E(\mathbf 1_E \mid \mathcal D_\alpha)$, the orthogonal projection of $\mathbf 1_E$ onto $L^2(\mathcal D_\alpha)$, is at least as close.  

\emph{A set in a single factor.}  Fix a Cantor set $\mathcal C$, integers $t \geq 1$ and $i \geq 0$, and $E \in \mathcal D$ with $m(E)>0$.  
Put $\kappa \coloneqq 1/(2(t+1))$.  Choose $\alpha$ such that $g \coloneqq \mathbb E(\mathbf 1_E \mid \mathcal D_\alpha)$ satisfies 
$\lVert \mathbf 1_E - g\rVert_2 ^2 < \kappa ^2 m(E)$, and let $H \coloneqq \{ g > 1-\kappa\} \in \mathcal D_\alpha$.  
Then $m(H)>0$: otherwise $1-g \geq \kappa$ almost everywhere on $E$, and so $\lVert \mathbf 1_E - g\rVert_2 ^2 \geq \kappa ^2 m(E)$.  

\emph{Comparison, for every step.}  Let $d$ be an integer, and $x \in H_d \coloneqq H \cap T ^{-d}H \cap \cdots \cap T ^{-td}H$.  
As $\mathcal D_\alpha$ is $T$-invariant, $\mathbb E(T ^{sd}\mathbf 1_E \mid \mathcal D_\alpha)(x) = g(T ^{sd}x) > 1-\kappa$ for $0 \leq s \leq t$.  
Since indicators take only the values $0$ and $1$, we have $\prod_{s=0}^t T ^{sd}\mathbf 1_E \geq 1 - \sum_{s=0}^t (1 - T ^{sd}\mathbf 1_E)$ pointwise, and so, on $H_d$, 
\begin{equation}
\mathbb E\Bigl( \prod_{s=0}^t T ^{sd}\mathbf 1_E \;\Big\vert\; \mathcal D_\alpha\Bigr) > 1 - (t+1) \kappa = \frac 12 . 
\end{equation}
As $H_d \in \mathcal D_\alpha$ and the product is non-negative, it follows that 
\begin{equation} \label{e:chainComparison}
m \bigl( E \cap T ^{-d}E \cap \cdots \cap T ^{-td}E \bigr) 
\geq \int_{H_d} \mathbb E\Bigl( \prod_{s=0}^t T ^{sd}\mathbf 1_E \;\Big\vert\; \mathcal D_\alpha\Bigr) dm 
\geq \frac 12 m(H_d) , \qquad d \in \mathbb Z . 
\end{equation}

\emph{Conclusion.}  Take $d = b ^{i}k_n$ in \eqref{e:chainComparison}, and average over $M \leq n < N$.  Since \eqref{e:chainComparison} holds for every $d$, 
\begin{equation}
\liminf_{N-M\to\infty} \Expectation_{M\leq n<N} m \Bigl( \bigcap_{s=0}^t T ^{-sb ^{i}k_n}E \Bigr) 
\geq \frac 12 \liminf_{N-M\to\infty} \Expectation_{M\leq n<N} m \Bigl( \bigcap_{s=0}^t T ^{-sb ^{i}k_n}H \Bigr) > 0 , 
\end{equation}
the last inequality because $H \in \mathcal D_\alpha$, $m(H)>0$, and $\mathcal D_\alpha$ is CSZ.  As $\mathcal C$, $t$, $i$ and $E$ were arbitrary, $\mathcal D$ is CSZ.  
\end{proof}

The comparison \eqref{e:chainComparison} is the one point at which the CSZ property enters, and it holds for every step $d$ separately.  
This is why the classical argument applies verbatim to the CSZ property, with its uniform averages over intervals of indices, 
its quantification over all Cantor sets, and its dilations: each of these is a statement about which steps $d$ are averaged.  

\begin{lemma}[Dichotomy] \label{l:dichotomy}
Let $(X,\mathcal A,m,T)$ be an ergodic system satisfying the standing assumption, and $\mathcal D \neq \mathcal A$ a factor.  
Then either $\mathcal A$ is a weakly mixing extension of $\mathcal D$, or there is a factor $\mathcal E$ with $\mathcal D \subsetneq \mathcal E \subset \mathcal A$ 
which is a compact extension of $\mathcal D$.  
\end{lemma}

This is due to Furstenberg \cite{MR0498471}; see \cite{MoreiraNotes}*{Proposition 8.6.1} for a proof in the present language.  
Our definitions differ slightly from those of \cite{MoreiraNotes}, and in the direction that makes the dichotomy apply as stated.  
There, $f$ is conditionally weakly mixing if $\lvert \langle T^nf,g\rangle_{\mathcal E\mid \mathcal D}\rvert \to 0$ in uniform Ces\`aro mean, in $L^2(\mathcal D)$, 
for \emph{every} $g$ in the Hilbert module \cite{MoreiraNotes}*{Definition 8.2.5}.  Our condition is the case $g=f$, with ordinary Ces\`aro means, and so is implied by it.  
Hence a weakly mixing extension in the sense of \cite{MoreiraNotes} is one in ours.  
Our definition of conditional compactness, with orbits indexed by $n \in \mathbb Z$ and approximants in the Hilbert module, is that of \cite{MoreiraNotes}*{Definition 8.2.8}.  
The proof in \cite{MoreiraNotes} is for ergodic systems on compact metric spaces, which is the setting of Theorem \ref{t:ECSZ}, and uses the disintegration of $m$ over $\mathcal D$.  

\begin{proof}[Proof of Theorem \ref{t:ECSZ}]
Let $\mathbf K$ be the set of CSZ factors of $(X,\mathcal A,m,T)$, partially ordered by inclusion.  
It is non-empty, as it contains the trivial factor; and every non-empty chain in $\mathbf K$ has an upper bound in $\mathbf K$, namely the $\sigma$-algebra generated by its union, by Lemma \ref{l:chain}.  
By Zorn's lemma, $\mathbf K$ has a maximal element $\mathcal D$.  

Suppose that $\mathcal D \neq \mathcal A$, and apply Lemma \ref{l:dichotomy}.  
If $\mathcal A$ is a weakly mixing extension of $\mathcal D$, then $\mathcal A$ is CSZ by Lemma \ref{l:Extensions}(W).  
Otherwise there is a compact extension $\mathcal E$ of $\mathcal D$ with $\mathcal D \subsetneq \mathcal E$, and $\mathcal E$ is CSZ by Lemma \ref{l:Extensions}(K).  
Either way, $\mathbf K$ contains a factor strictly larger than $\mathcal D$, contradicting maximality.  
Hence $\mathcal D = \mathcal A$, which is to say that the system is CSZ.  
\end{proof}


Finally, we record the van der Corput inequality, which is the starting point of the argument in \S\ref{sec:wm}. 

\begin{lemma}[van der Corput Inequality]  \label{vdC}
Let $u_M, \ldots, u_{N-1}$ be vectors in a Hilbert space, of norm at most $1$, and let $1\leq H < N-M$.  Then 
\begin{equation}
    \label{vdc} 
\Bigl\lVert \Expectation_{M\leq n <N} u_n \Bigr\rVert ^2 
\leq 
4\, \Expectation_{1\leq h\leq H} 
\Bigl\lvert 
\Expectation_{M\leq n <N-h} \langle u_{n+h}, u_n\rangle  
\Bigr\rvert
+ \frac 2H . 
\end{equation}
\end{lemma}

\begin{proof}
Write $L = N-M$, and set $u_n = 0$ for $n \notin [M,N)$.  Each $u_n$ with $n \in [M,N)$ occurs exactly $H$ times among the vectors $u_{m+h}$, 
$1 \leq h \leq H$, $m \in I \coloneqq [M-H, N-1)$, an interval of fewer than $L+H$ integers.  By Cauchy--Schwarz in $m$, 
\begin{equation}
\Bigl\lVert \sum_{n} u_n \Bigr\rVert ^2 
= \Bigl\lVert \sum _{m \in I} \frac 1H \sum _{h=1}^H u_{m+h} \Bigr\rVert ^2 
\leq \frac{L+H}{H^2} \sum _{h,h'=1}^H \sum _{m \in I} \langle u_{m+h}, u_{m+h'} \rangle . 
\end{equation}
The $H$ diagonal terms contribute at most $L$ each.  For $r = h-h' \neq 0$, the inner sum is 
$S(r) \coloneqq \sum _{n} \langle u_{n+r}, u_n\rangle$, with $S(-r) = \overline{S(r)}$, and exactly $H - \lvert r\rvert$ pairs $(h,h')$ have $h-h'=r$.  
Moreover $\lvert S(r)\rvert = (L-r) \lvert c_r\rvert \leq L \lvert c_r\rvert$ for $1 \leq r <H$, where $c_r \coloneqq \Expectation_{M\leq n<N-r}\langle u_{n+r},u_n\rangle$.  Hence 
\begin{equation}
\Bigl\lVert \Expectation_{M\leq n<N} u_n \Bigr\rVert ^2 
\leq \frac{L+H}{HL} + \frac{2(L+H)}{H^2 L} \sum _{r=1}^{H-1} (H-r) \lvert c_r\rvert 
\leq \frac 1H + \frac 1L + \frac{2(L+H)}{L}\, \Expectation_{1\leq r\leq H} \lvert c_r\rvert . 
\end{equation}
As $H < L$, this is \eqref{vdc}. 
\end{proof}


\subsection{Reduction to ergodic systems on compact metric spaces} \label{sub:reduction}

The CSZ property asks for a positive liminf over all intervals of indices, which is more than the usual multiple recurrence statement, 
so we give the reduction to ergodic systems in full.  

\begin{proposition} \label{p:reduction}
Suppose that every ergodic system satisfying the standing assumption of \S\ref{sub:mps} is CSZ.  Then every invertible measure preserving system is CSZ.  
\end{proposition}

\begin{proof}
Fix an invertible measure preserving system $(X,\mathcal A,m,T)$, a Cantor set $\mathcal C$ of base $b$, integers $t \geq 1$ and $i \geq 0$, and $A \in \mathcal A$ with $m(A)>0$.  

\emph{A compact metric model.}  Let $\Omega \coloneqq \{0,1\} ^{\mathbb Z}$, a compact metric space in the product topology, with the shift $(\tau\omega)(n) \coloneqq \omega(n+1)$, a homeomorphism.  
Define $\Phi \colon X \to \Omega$ by $\Phi(x)(n) \coloneqq \mathbf 1_A(T ^{n}x)$.  Each coordinate of $\Phi$ is measurable, so $\Phi$ is measurable, and $\Phi\circ T = \tau \circ \Phi$.  
Hence $\mu \coloneqq \Phi_\ast m$ is a $\tau$-invariant Borel probability measure on $\Omega$.  Let $\widetilde A \coloneqq \{\omega \colon \omega(0) = 1\}$.  
Then $\Phi ^{-1}(\tau ^{-d}\widetilde A) = T ^{-d}A$ for every $d \in \mathbb Z$, and therefore 
\begin{equation} \label{e:model}
m \Bigl( \bigcap_{s=0}^t T ^{-d_s}A \Bigr) = \mu \Bigl( \bigcap_{s=0}^t \tau ^{-d_s}\widetilde A \Bigr) , \qquad d_0 , \ldots, d_t \in \mathbb Z . 
\end{equation}
In particular $\mu(\widetilde A) = m(A)>0$.  By \eqref{e:model}, the averages in \eqref{e:CSZdef} for $A$ in $(X,\mathcal A,m,T)$ coincide with those for $\widetilde A$ in $(\Omega,\mathcal B,\mu,\tau)$, 
with $\mathcal B$ the Borel $\sigma$-algebra of $\Omega$.  So it suffices to prove that the latter have a positive liminf.  

\emph{Ergodic decomposition.}  By \cite{MoreiraNotes}*{Theorem 7.3.1}, for $\mu$-almost every $y \in \Omega$ there is an ergodic $\tau$-invariant Borel probability measure $\mu_y$ on $\Omega$ 
such that, for every Borel set $B$, the map $y \mapsto \mu_y(B)$ is measurable and $\mu(B) = \int \mu_y(B)\;d\mu(y)$.  
For a $\tau$-invariant Borel probability measure $\nu$, write 
$c_n(\nu) \coloneqq \nu \bigl( \bigcap_{s=0}^t \tau ^{-sb ^{i}k_n}\widetilde A \bigr)$, so that $c_n(\mu) = \int c_n(\mu_y)\;d\mu(y)$ for every $n$.  
As $0 < \mu(\widetilde A) = \int \mu_y(\widetilde A)\;d\mu(y)$, the set $Y_+ \coloneqq \{y \colon \mu_y(\widetilde A)>0\}$ has positive measure.  
For $y \in Y_+$, the system $(\Omega,\mathcal B,\mu_y,\tau)$ is ergodic, invertible, and on a compact metric space, so it is CSZ by hypothesis, and 
\begin{equation}
c(y) \coloneqq \liminf_{N-M\to\infty} \Expectation_{M\leq n<N} c_n(\mu_y) > 0 , \qquad y \in Y_+ . 
\end{equation}

\emph{Passing the liminf through the integral.}  Let $c \coloneqq \liminf_{N-M\to\infty} \Expectation_{M\leq n<N} c_n(\mu)$, 
and choose intervals $[M_k,N_k)$, with $N_k - M_k \to \infty$, along which $\Expectation_{M_k\leq n<N_k} c_n(\mu) \to c$.  
The function $\ell(y) \coloneqq \liminf_k \Expectation_{M_k\leq n<N_k} c_n(\mu_y)$ is measurable, and $\ell(y) \geq c(y)$, 
because a liminf along one sequence of intervals whose lengths tend to infinity is at least the liminf over all such intervals.  
By Fatou's lemma, the integrands being non-negative, 
\begin{equation}
c = \lim_k \int \Expectation_{M_k\leq n<N_k} c_n(\mu_y)\;d\mu(y) \geq \int \ell(y)\;d\mu(y) \geq \int_{Y_+} \ell(y)\;d\mu(y) > 0 , 
\end{equation}
the last inequality because $\ell(y) \geq c(y) > 0$ for $y \in Y_+$, a set of positive measure.  As $\mathcal C$, $t$, $i$ and $A$ were arbitrary, $(X,\mathcal A,m,T)$ is CSZ.  
\end{proof}

The step through Fatou's lemma is the only point at which the uniformity of the CSZ property makes a difference here.  
A uniform liminf is not a liminf along a single sequence, so Fatou's lemma cannot be applied to it directly; 
it is applied instead along one sequence of intervals that realizes the liminf for $\mu$.  
Compare \cite{MoreiraNotes}*{\S 7.5}, where the corresponding reduction is carried out for uniform limits.  

We record where the hypotheses of the standing assumption, and ergodicity, are used.  
The compact metric structure is used for three things: the ergodic decomposition in Proposition \ref{p:reduction}; 
the disintegration of $m$ over a factor, behind the joining $\lambda$ of Lemma \ref{l:atomless} and so behind \S\ref{sec:wm}; 
and the dichotomy, Lemma \ref{l:dichotomy}, whose proof in \cite{MoreiraNotes} uses the same disintegration.  
The approximation step in the proof of Lemma \ref{l:chain} does not use it.  
Ergodicity is used only in the dichotomy.  Lemmas \ref{l:Extensions} and \ref{l:chain}, and Theorems \ref{t:compactSystem} and \ref{t:compactExtension}, hold without it.  
